\documentclass[../../main.tex]{subfiles}% 注意这里的文件路径不能用 ./main.tex ,否则用latexmk编译子文件会报错
\graphicspath{{\subfix{./image/}}} % 指定图片目录,后续可以直接使用图片文件名
% 注意这里的文件路径不能用 ../../image/ ,否则用latexmk编译子文件会报错

% 例如:
% \begin{figure}[H]
% \centering
% \includegraphics[scale=0.3]{图.png}
% \caption{}
% \label{figure:图}
% \end{figure}
% 注意:上述\label{}一定要放在\caption{}之后,否则引用图片序号会只会显示??.

\begin{document}

\section{曲面论基本方程}

\begin{definition}
设参数区域$D\subseteq\mathbb{R}^2$，在$D$上给定两个对称二次微分形式
\begin{align}
\varphi &= g_{\alpha\beta}\mathrm{d}u^\alpha\mathrm{d}u^\beta,\quad \psi = b_{\alpha\beta}\mathrm{d}u^\alpha\mathrm{d}u^\beta, \label{eq:form_def}
\end{align}
其中$\alpha,\beta=1,2$，$g_{\alpha\beta}=g_{\beta\alpha}$，$b_{\alpha\beta}=b_{\beta\alpha}$，且$\varphi$为正定二次形式。若存在正则参数曲面$\boldsymbol{f}:D\to\mathbb{E}^3$，使得$\varphi,\psi$分别为$\boldsymbol{f}$的第一、第二基本形式，则称$\varphi,\psi$可实现为欧氏空间曲面的基本形式。
\end{definition}

\begin{definition}
设正则参数曲面$\boldsymbol{r}=\boldsymbol{r}(u^1,u^2)$的第一基本形式系数为$g_{\alpha\beta}$，定义\textbf{第一类}$\mathbf{Christoffel}$\textbf{记号}
\begin{align}
\Gamma^\gamma_{\alpha\beta} = \frac{1}{2}g^{\gamma\xi}\left(\frac{\partial g_{\alpha\xi}}{\partial u^\beta}+\frac{\partial g_{\xi\beta}}{\partial u^\alpha}-\frac{\partial g_{\alpha\beta}}{\partial u^\xi}\right), \label{eq:christoffel}
\end{align}
其中$(g^{\alpha\beta})=(g_{\alpha\beta})^{-1}$为度量矩阵的逆矩阵；定义指标提升形式
\begin{align*}
b^\gamma_\beta = g^{\gamma\xi}b_{\xi\beta},
\end{align*}
其中$b_{\alpha\beta}$为曲面第二基本形式系数。
\end{definition}

\begin{proposition}
正则参数曲面$\boldsymbol{r}=\boldsymbol{r}(u^1,u^2)$的自然标架场$\{\boldsymbol{r};\boldsymbol{r}_1,\boldsymbol{r}_2,\boldsymbol{n}\}$满足\textbf{自然标架运动公式}
\begin{align}
\begin{cases}
\frac{\partial \boldsymbol{r}}{\partial u^\alpha} = \boldsymbol{r}_\alpha,\\
\frac{\partial \boldsymbol{r}_\alpha}{\partial u^\beta} = \Gamma^\gamma_{\alpha\beta}\boldsymbol{r}_\gamma + b_{\alpha\beta}\boldsymbol{n},\\
\frac{\partial \boldsymbol{n}}{\partial u^\beta} = -b^\gamma_\beta\boldsymbol{r}_\gamma.
\end{cases} \label{eq:motion_formula111}
\end{align}
\end{proposition}

\begin{definition}
由第一类基本量$g_{\alpha\beta}$及其偏导数定义$\mathbf{Riemann}$\textbf{记号}
\begin{align}
R^\delta_{\alpha\beta\gamma} = \frac{\partial \Gamma^\delta_{\alpha\beta}}{\partial u^\gamma}-\frac{\partial \Gamma^\delta_{\alpha\gamma}}{\partial u^\beta}+\Gamma^\eta_{\alpha\beta}\Gamma^\delta_{\eta\gamma}-\Gamma^\eta_{\alpha\gamma}\Gamma^\delta_{\eta\beta}, \label{eq:riemann_up}
\end{align}
其协变形式与逆变形式的升降关系为
\begin{align}
R_{\alpha\delta\beta\gamma} = g_{\delta\eta}R^\eta_{\alpha\beta\gamma},\quad R^\delta_{\alpha\beta\gamma} = g^{\delta\eta}R_{\alpha\eta\beta\gamma}. \label{eq:riemann_cov}
\end{align}
\end{definition}

\begin{theorem}[Gauss-Codazzi相容性方程]
正则参数曲面的第一、第二基本形式系数必满足
\begin{gather}
R^\delta_{\alpha\beta\gamma} = b_{\alpha\beta}b^\delta_\gamma - b_{\alpha\gamma}b^\delta_\beta, \tag{Gauss方程}\label{eq:gauss_eq}\\
\frac{\partial b_{\alpha\beta}}{\partial u^\gamma}-\frac{\partial b_{\alpha\gamma}}{\partial u^\beta} = \Gamma^\delta_{\alpha\gamma}b_{\delta\beta}-\Gamma^\delta_{\alpha\beta}b_{\delta\gamma}. \tag{Codazzi方程}\label{eq:codazzi_eq}
\end{gather}
\end{theorem}

\begin{proof}
正则参数曲面具有三阶以上连续偏导数，故$\boldsymbol{r}_\alpha,\boldsymbol{n}$的二阶混合偏导与求导次序无关，即
\begin{align*}
\frac{\partial^2 \boldsymbol{r}_\alpha}{\partial u^\beta\partial u^\gamma} = \frac{\partial^2 \boldsymbol{r}_\alpha}{\partial u^\gamma\partial u^\beta},\quad \frac{\partial^2 \boldsymbol{n}}{\partial u^\beta\partial u^\gamma} = \frac{\partial^2 \boldsymbol{n}}{\partial u^\gamma\partial u^\beta}.
\end{align*}
将运动公式\eqref{eq:motion_formula111}代入$\boldsymbol{r}_\alpha$的二阶偏导等式：
\begin{align*}
\frac{\partial}{\partial u^\gamma}\left(\Gamma^\delta_{\alpha\beta}\boldsymbol{r}_\delta + b_{\alpha\beta}\boldsymbol{n}\right) = \frac{\partial}{\partial u^\beta}\left(\Gamma^\delta_{\alpha\gamma}\boldsymbol{r}_\delta + b_{\alpha\gamma}\boldsymbol{n}\right).
\end{align*}
展开并再次代入运动公式，整理得
\begin{align*}
&\left(\frac{\partial \Gamma^\delta_{\alpha\beta}}{\partial u^\gamma}-\frac{\partial \Gamma^\delta_{\alpha\gamma}}{\partial u^\beta}+\Gamma^\eta_{\alpha\beta}\Gamma^\delta_{\eta\gamma}-\Gamma^\eta_{\alpha\gamma}\Gamma^\delta_{\eta\beta}-b_{\alpha\beta}b^\delta_\gamma + b_{\alpha\gamma}b^\delta_\beta\right)\boldsymbol{r}_\delta\\
&+\left(\Gamma^\delta_{\alpha\beta}b_{\delta\gamma}-\Gamma^\delta_{\alpha\gamma}b_{\delta\beta}+\frac{\partial b_{\alpha\beta}}{\partial u^\gamma}-\frac{\partial b_{\alpha\gamma}}{\partial u^\beta}\right)\boldsymbol{n} = \boldsymbol{0}.
\end{align*}
由于$\boldsymbol{r}_1,\boldsymbol{r}_2,\boldsymbol{n}$线性无关，其系数分别为零，即得\eqref{eq:gauss_eq}与\eqref{eq:codazzi_eq}。

对$\boldsymbol{n}$的二阶偏导等式代入运动公式，展开得
\begin{align*}
\frac{\partial b^\delta_\beta}{\partial u^\gamma}-\frac{\partial b^\delta_\gamma}{\partial u^\beta} = -b^\eta_\beta\Gamma^\delta_{\eta\gamma}+b^\eta_\gamma\Gamma^\delta_{\eta\beta},
\end{align*}
该式与Codazzi方程\eqref{eq:codazzi_eq}等价，无新相容性条件。

\end{proof}

\begin{remark}
Riemann记号具有对称性质
\begin{align*}
R_{\alpha\delta\beta\gamma} = R_{\beta\gamma\alpha\delta} = -R_{\delta\alpha\beta\gamma} = -R_{\alpha\delta\gamma\beta},
\end{align*}
由此可得$R_{11\beta\gamma}=R_{22\beta\gamma}=0$，$R_{\alpha\delta11}=R_{\alpha\delta22}=0$。因此：
\begin{enumerate}[(1)]
\item Gauss方程仅有1个独立等式：
\begin{align}
R_{1212} = b_{11}b_{22}-(b_{12})^2; \label{eq:gauss_indep}
\end{align}
\item Codazzi方程仅有2个独立等式（$\beta=1,\gamma=2,\alpha=1,2$）：
\begin{align}
\begin{cases}
\frac{\partial b_{11}}{\partial u^2}-\frac{\partial b_{12}}{\partial u^1} = -b_{2\delta}\Gamma^\delta_{11}+b_{1\delta}\Gamma^\delta_{12},\\
\frac{\partial b_{21}}{\partial u^2}-\frac{\partial b_{22}}{\partial u^1} = -b_{2\delta}\Gamma^\delta_{21}+b_{1\delta}\Gamma^\delta_{22}.
\end{cases} \label{eq:codazzi_indep}
\end{align}
\end{enumerate}
\end{remark}

\begin{theorem}[曲面基本形式存在唯一性定理]
设参数区域$D$上给定对称正定二次微分形式$\varphi=g_{\alpha\beta}\mathrm{d}u^\alpha\mathrm{d}u^\beta$与对称二次微分形式$\psi=b_{\alpha\beta}\mathrm{d}u^\alpha\mathrm{d}u^\beta$，若其系数满足Gauss-Codazzi方程\eqref{eq:gauss_eq},\eqref{eq:codazzi_eq}，则存在正则参数曲面$\boldsymbol{f}:D\to\mathbb{E}^3$，使得$\varphi,\psi$分别为其第一、第二基本形式，且在给定初始条件下曲面唯一。
\end{theorem}
\begin{remark}
曲线论中曲率、挠率可取弧长的任意正函数；曲面论中第一、第二基本形式不独立，保持第一基本形式不变的变形必保持曲面内蕴弯曲性质不变，此为$\mathbf{Gauss}$\textbf{绝妙定理}的核心内涵。
\end{remark}
\begin{proof}
Gauss-Codazzi方程等价于运动公式\eqref{eq:motion_formula111}的相容性条件（二阶混合偏导可交换）。运动公式为含12个数量未知函数的一阶偏微分方程组，由一阶偏微分方程组解的存在唯一性定理，该方程组可积，即解存在且唯一，对应所求曲面。

\end{proof}

\begin{proposition}
设曲面取正交参数曲线网$(u,v)$（即第一基本形式中$F=0$，$\mathrm{I}=E\mathrm{d}u^2+G\mathrm{d}v^2$），则
\begin{align*}
R_{1212} = -\sqrt{EG}\left(\left(\frac{(\sqrt{E})_v}{\sqrt{G}}\right)_v+\left(\frac{(\sqrt{G})_u}{\sqrt{E}}\right)_u\right).
\end{align*}
\end{proposition}

\begin{proof}
正交参数下Christoffel记号有简化形式，代入Riemann记号定义\eqref{eq:riemann_up}，整理化简即得。

\end{proof}

\begin{proposition}
设曲面取正交曲率线网为参数曲线网（即$F=M=0$），记平均曲率$H=\frac{1}{2}\left(\frac{L}{E}+\frac{N}{G}\right)$，则Codazzi方程简化为
\begin{align*}
\frac{\partial L}{\partial v}=H\frac{\partial E}{\partial v},\quad \frac{\partial N}{\partial u}=H\frac{\partial G}{\partial u}.
\end{align*}
\end{proposition}
\begin{proof}
曲率线参数下$F=M=0$，将正交参数下Christoffel记号代入Codazzi独立方程\eqref{eq:codazzi_indep}，化简即得。

\end{proof}

\begin{example}
验证方程组$(3.14)$和$(3.9)$的等价性。
\end{example}

\begin{example}
证明: 若$(u,v)$是曲面$S$上的参数系, 使得参数曲线网是曲面$S$上的正交曲率线网, 则主曲率$\kappa_1,\kappa_2$满足方程组
\begin{align*}
\frac{\partial \kappa_1}{\partial v} = \frac{E_v}{2E}(\kappa_2 - \kappa_1),\quad \frac{\partial \kappa_2}{\partial u} = \frac{G_u}{2G}(\kappa_1 - \kappa_2).
\end{align*}
\end{example}

\begin{example}
证明: 平均曲率为常数的曲面或者是平面, 或者是球面, 或者存在参数系$(u,v)$, 使得它的第一基本形式和第二基本形式能够表示成
\begin{align*}
\mathrm{I} &= \lambda\left((\mathrm{d}u)^2 + (\mathrm{d}v)^2\right),\\
\mathrm{II} &= (1+\lambda H)(\mathrm{d}u)^2 - (1-\lambda H)(\mathrm{d}v)^2,
\end{align*}
其中$H$是曲面的平均曲率。
\end{example}

\begin{example}
设$S$是$E^3$中的一块曲面, 它的主曲率是两个互不相等的常值函数, 证明: $S$是圆柱面的一部分。
\end{example}

\begin{example}
已知曲面$S$的第一基本形式和第二基本形式分别是
\begin{align*}
\mathrm{I} &= u^2\left((\mathrm{d}u)^2 + (\mathrm{d}v)^2\right),\\
\mathrm{II} &= A(u,v)(\mathrm{d}u)^2 + B(u,v)(\mathrm{d}v)^2,
\end{align*}
证明: 函数$A(u,v), B(u,v)$满足关系式
\begin{align*}
A(u,v)\cdot B(u,v) = 1,
\end{align*}
并且它们只是$u$的函数。
\end{example}





\end{document}